% =============================================================================
%  PAES Studio: guía de síntesis (plantilla)
%  Lista para Overleaf: Nuevo proyecto → Subir proyecto → compilar con pdfLaTeX.
%
%  Formato: A4, dos columnas, cajas tcolorbox con la paleta Zinc de la app.
%  Entornos disponibles:
%    definicion, formula, ley, mecanismo, ideasclave, metodo, ejemplo,
%    errorfrecuente, dato, lineatiempo   (título opcional en [])
%  Comandos:
%    \alternativa{A}{texto}   alternativa con su tecla (igual que en la app)
%    \clave{C}        marca la alternativa correcta
%
%  El generador scripts/pipeline/generate_latex_summary.py reemplaza los bloques
%  delimitados por %%<NOMBRE>%% … %%</NOMBRE>%%. Fuera de ellos, edita libremente.
% =============================================================================
\documentclass[10pt,a4paper,twocolumn]{article}

\usepackage{lmodern}        % antes de fontenc/babel: no depende de las fuentes EC
\usepackage[T1]{fontenc}
\usepackage[utf8]{inputenc}
\usepackage[spanish,es-noshorthands,es-tabla]{babel}
\spanishplainpercent        % \% sin espaciado automático (evita choques en modo matemático)
\usepackage{textcomp}
\usepackage{amsmath,amssymb}
\usepackage[a4paper,margin=1.5cm,top=1.9cm,bottom=1.9cm,columnsep=0.75cm]{geometry}
\usepackage{xcolor}
\usepackage{enumitem}
\usepackage{booktabs}
\usepackage{array}
\usepackage{titlesec}
\usepackage{fancyhdr}
\usepackage[most]{tcolorbox}
\usepackage[hidelinks]{hyperref}

% ── Paleta Zinc (misma que la app) ───────────────────────────────────────────
\definecolor{zinc950}{HTML}{09090B}
\definecolor{zinc900}{HTML}{18181B}
\definecolor{zinc700}{HTML}{3F3F46}
\definecolor{zinc500}{HTML}{71717A}
\definecolor{zinc300}{HTML}{D4D4D8}
\definecolor{zinc200}{HTML}{E4E4E7}
\definecolor{zinc100}{HTML}{F4F4F5}
\definecolor{zinc50}{HTML}{FAFAFA}
\definecolor{esmeralda}{HTML}{047857}
\definecolor{rosa}{HTML}{BE123C}
\definecolor{ambar}{HTML}{B45309}

% ── Tipografía: texto con serifa, estructura en sans ─────────────────────────
\setlength{\parindent}{0pt}
\setlength{\parskip}{0.45em}
\linespread{1.08}
\color{zinc900}
\newcommand{\sen}{\operatorname{sen}}
\renewcommand{\arraystretch}{1.18}

\titleformat{\section}{\sffamily\bfseries\large\color{zinc950}}{\thesection}{0.6em}{}
  [{\color{zinc300}\titlerule[0.4pt]}]
\titleformat{\subsection}{\sffamily\bfseries\normalsize\color{zinc900}}{\thesubsection}{0.5em}{}
\titlespacing*{\section}{0pt}{1.1em}{0.5em}
\titlespacing*{\subsection}{0pt}{0.9em}{0.3em}

\setlist{leftmargin=1.2em,itemsep=0.15em,topsep=0.3em}
\setlist[itemize]{label={\color{zinc500}\textbullet}}

\pagestyle{fancy}
\fancyhf{}
\renewcommand{\headrulewidth}{0pt}
\fancyfoot[L]{\sffamily\scriptsize\color{zinc500}\guiaAsignatura}
\fancyfoot[R]{\sffamily\scriptsize\color{zinc500}\thepage}

% ── Cajas ────────────────────────────────────────────────────────────────────
\tcbset{
  paes/.style={
    enhanced, breakable, sharp corners=south, arc=1.5pt,
    boxrule=0.4pt, colback=white, colframe=zinc300,
    colbacktitle=white, coltitle=zinc900, titlerule=0pt,
    fonttitle=\sffamily\bfseries\footnotesize,
    toptitle=3pt, bottomtitle=0pt,
    left=7pt, right=7pt, top=2pt, bottom=6pt,
    before skip=7pt, after skip=7pt,
  },
  acento/.style={borderline west={2pt}{0pt}{#1}},
}
\newtcolorbox{definicion}[1][Definición]{paes, acento=zinc900, title={#1}}
\newtcolorbox{formula}[1][Fórmula clave]{paes, colback=zinc50, colbacktitle=zinc50, title={#1}}
\newtcolorbox{ejemplo}[1][Ejemplo resuelto]{paes, title={#1}}
\newtcolorbox{errorfrecuente}[1][Error frecuente]{paes, colframe=rosa!35, acento=rosa, title={#1}}
\newtcolorbox{dato}[1][Dato PAES]{paes, colframe=ambar!40, acento=ambar, title={#1}}
\newtcolorbox{mecanismo}[1][Mecanismo]{paes, colback=zinc50, colbacktitle=zinc50, acento=zinc700, title={#1}}
\newtcolorbox{ley}[1][Ley]{paes, colback=zinc50, colbacktitle=zinc50, acento=zinc900, title={#1}}
\newtcolorbox{ideasclave}[1][Ideas clave]{paes, title={#1}}
\newtcolorbox{metodo}[1][Cómo abordarlo en la PAES]{paes, acento=esmeralda, title={#1}}
% Línea de tiempo: filas «fecha & hecho \\».
\newenvironment{lineatiempo}[1][Línea de tiempo]
  {\begin{tcolorbox}[paes, acento=zinc700, title={#1}]\small\renewcommand{\arraystretch}{1.25}%
   \begin{tabular}{@{}>{\leavevmode\sffamily\bfseries\color{zinc700}}p{0.24\linewidth}p{\dimexpr0.76\linewidth-2\tabcolsep\relax}@{}}}
  {\end{tabular}\end{tcolorbox}}

% Alternativa con su tecla (eco de la anatomía A–E de la app).
\newtcbox{\tecla}{on line, boxrule=0.4pt, arc=1.5pt, colback=zinc50, colframe=zinc300,
  left=2.2pt, right=2.2pt, top=0.6pt, bottom=0.6pt, fontupper=\ttfamily\scriptsize\color{zinc700}}
\newtcbox{\teclaok}{on line, boxrule=0.4pt, arc=1.5pt, colback=esmeralda, colframe=esmeralda,
  left=2.2pt, right=2.2pt, top=0.6pt, bottom=0.6pt, fontupper=\ttfamily\scriptsize\bfseries\color{white}}
\newcommand{\alternativa}[2]{\par\noindent\hangindent=1.6em\hangafter=1\makebox[1.6em][l]{\tecla{#1}}#2\par}
\newcommand{\clave}[1]{\par\smallskip{\sffamily\footnotesize\color{esmeralda} Clave \teclaok{#1}}\par}

% Barra horizontal proporcional (para tablas de frecuencia).
\newcommand{\barra}[1]{{\color{zinc900}\rule[0.15ex]{#1\linewidth}{0.9ex}}}

%%<META>%%
\newcommand{\guiaTitulo}{Compendio de estudio}
\newcommand{\guiaAsignatura}{Ciencias}
\newcommand{\guiaDescripcion}{La prueba de Ciencias tiene 80 preguntas: 54 del módulo común (18 de Biología, 18 de Física y 18 de Química) y 26 del módulo electivo. Además de contenidos, evalúa habilidades de pensamiento científico: identificar variables, formular hipótesis, interpretar datos y evaluar conclusiones.}
\newcommand{\guiaFuente}{4 ejes y 17 subtemas. Frecuencias calculadas sobre 21 ensayos (20 oficiales DEMRE), 1474 preguntas. Generado el 8 de noviembre de 2026.}
%%</META>%%

\begin{document}

\twocolumn[{%
  \begin{tcolorbox}[enhanced, frame hidden, colback=white, left=0pt, right=0pt, top=0pt, bottom=6pt,
      borderline south={0.4pt}{0pt}{zinc300}]
    {\sffamily\footnotesize\color{zinc500} PAES Studio}\par\vspace{0.2em}
    {\sffamily\bfseries\LARGE\color{zinc950} \guiaTitulo: \guiaAsignatura}\par\vspace{0.35em}
    {\normalsize\color{zinc700} \guiaDescripcion}\par\vspace{0.25em}
    {\sffamily\scriptsize\color{zinc500} \guiaFuente}
  \end{tcolorbox}
  \vspace{0.6em}
}]

%%<CONTENIDO>%%
\section*{Cómo usar este compendio}

Cada subtema sigue el mismo orden: teoría, recuadro con las fórmulas, leyes o mecanismos centrales, ideas clave, método para responder en la PAES, trampas frecuentes y ejemplos resueltos paso a paso. Resuelve cada ejemplo antes de leer su desarrollo.

\section*{Qué se evalúa}

Distribución de las preguntas por eje en los ensayos disponibles en PAES Studio.

\begin{center}\small

\begin{tabular}{@{}p{0.42\linewidth}rrp{0.22\linewidth}@{}}

\toprule Eje & Preg. & \% & \\ \midrule

Biología & 498 & 34 & \barra{0.074} \\

Física & 498 & 34 & \barra{0.074} \\

Química & 478 & 32 & \barra{0.071} \\

\bottomrule

\end{tabular}

\end{center}

\begin{dato}Biología y Física concentran el 68 \% de las preguntas: priorízalos si tienes poco tiempo.\end{dato}

\section{Biología}

{\itshape\color{zinc700} Organización celular, procesos y funciones biológicas (sistema nervioso y reproducción), herencia y evolución, organismo y ambiente.}

\subsection{Organización, estructura y actividad celular}

La célula es la unidad estructural y funcional de los seres vivos. Las procariontes (bacterias) no tienen núcleo ni organelos con membrana; las eucariontes (animales, plantas, hongos, protistas) tienen núcleo y organelos membranosos.

La membrana plasmática es una bicapa de fosfolípidos con proteínas que regula el paso de sustancias. El transporte pasivo (difusión, osmosis) ocurre a favor del gradiente y sin gasto de ATP; el activo, en contra del gradiente y con gasto de ATP.

\begin{mecanismo}[Organelos y función]
Núcleo: material genético. Mitocondria: respiración celular (ATP). Cloroplasto: fotosíntesis (solo vegetales). Ribosoma: síntesis de proteínas. Retículo endoplasmático y aparato de Golgi: síntesis, modificación y transporte. Lisosoma: digestión celular.
\end{mecanismo}

{\sffamily\footnotesize\color{zinc500} Esquema comparativo}
\begin{center}\small
\begin{tabular}{@{}>{\raggedright\arraybackslash}p{\dimexpr(\linewidth-4\tabcolsep)/3\relax}>{\raggedright\arraybackslash}p{\dimexpr(\linewidth-4\tabcolsep)/3\relax}>{\raggedright\arraybackslash}p{\dimexpr(\linewidth-4\tabcolsep)/3\relax}@{}}
\toprule
 & \textbf{Procarionte} & \textbf{Eucarionte} \\
\midrule
Núcleo & no (nucleoide) & sí \\
Organelos con membrana & no & sí \\
Ejemplos & bacterias & animales, plantas, hongos \\
\bottomrule
\end{tabular}
\end{center}

\begin{ideasclave}
\begin{itemize}
  \item Célula vegetal: pared celular, cloroplastos y vacuola grande.
  \item Osmosis: el agua pasa hacia el medio más concentrado en solutos.
  \item Los virus no son células.
\end{itemize}
\end{ideasclave}

\begin{metodo}
\begin{enumerate}
  \item Identifica el tipo de célula o estructura involucrada.
  \item Relaciona estructura con función.
  \item En transporte, determina el sentido del gradiente y si hay gasto de energía.
\end{enumerate}
\end{metodo}

\begin{errorfrecuente}[Trampas frecuentes del DEMRE]
\begin{itemize}
  \item Atribuir cloroplastos o pared celular a células animales.
  \item Pensar que en la osmosis se mueve el soluto (se mueve el agua).
\end{itemize}
\end{errorfrecuente}

\begin{ejemplo}[Ejemplo resuelto]
{\sffamily\footnotesize\color{zinc500} Enunciado}\par
Un glóbulo rojo se coloca en agua destilada. ¿Qué ocurrirá con la célula?
\alternativa{A}{Se deshidrata, porque el agua sale de ella.}
\alternativa{B}{Se hincha y puede estallar, porque el agua entra en ella.}
\alternativa{C}{No cambia, porque la membrana es impermeable al agua.}
\alternativa{D}{Pierde sus sales hacia el agua por transporte activo.}
\clave{B}
{\sffamily\footnotesize\color{zinc500} Análisis}\par El agua destilada es hipotónica respecto del interior celular.
{\sffamily\footnotesize\color{zinc500} Desarrollo}
\begin{enumerate}
  \item El interior del glóbulo rojo tiene más solutos que el agua destilada.
  \item Por osmosis, el agua entra hacia el medio más concentrado: la célula se hincha (hemólisis).
\end{enumerate}
\end{ejemplo}

\begin{ejemplo}[Ejemplo tipo PAES 2: Enzimas]
{\sffamily\footnotesize\color{zinc500} Enunciado}\par
Un grupo de estudiantes mide la actividad de una enzima digestiva humana a distintas temperaturas. Observan que la actividad aumenta hasta un máximo cerca de 37 \textdegree{}C y que sobre 45 \textdegree{}C cae bruscamente hasta casi desaparecer.

¿Qué explica la caída de la actividad sobre 45 \textdegree{}C?
\alternativa{A}{La enzima se consume a medida que cataliza la reacción.}
\alternativa{B}{La enzima pierde su estructura tridimensional y con ello su sitio activo funcional.}
\alternativa{C}{El sustrato se transforma en enzima a altas temperaturas.}
\alternativa{D}{Las moléculas se mueven más lento y chocan con menos frecuencia.}
\clave{B}
{\sffamily\footnotesize\color{zinc500} Análisis}\par El calor excesivo desnaturaliza la enzima: rompe las interacciones que mantienen su forma, el sitio activo deja de encajar con el sustrato y la actividad cae.
{\sffamily\footnotesize\color{zinc500} Desarrollo}
\begin{enumerate}
  \item La enzima es una proteína y su función depende de su estructura tridimensional, que se mantiene gracias a enlaces débiles como los puentes de hidrógeno y las interacciones iónicas.
  \item Hasta la temperatura óptima, cerca de 37 \textdegree{}C, el calor aumenta la energía cinética de las moléculas y la frecuencia de los choques entre enzima y sustrato; por eso la actividad sube.
  \item Sobre 45 \textdegree{}C la energía térmica supera la estabilidad de esos enlaces y la proteína se desnaturaliza: se despliega y cambia su conformación.
  \item Al cambiar su forma, el sitio activo deja de ser complementario al sustrato. Los sustratos ya no encajan y la catálisis cae hasta casi desaparecer.
  \item La alternativa A es falsa: las enzimas son catalizadores y no se consumen en la reacción, así que su caída no se explica por consumo.
  \item La alternativa C no tiene sustento: el sustrato no se convierte en enzima, pues son moléculas distintas, y el sustrato se transforma en producto.
  \item La alternativa D describe lo contrario de lo ocurrido: a mayor temperatura las moléculas se mueven más rápido, lo que explica el aumento inicial y no la caída.
  \item Así, la explicación de la caída es la pérdida de la estructura tridimensional y, con ella, del sitio activo. Por lo tanto, la alternativa correcta es la B.
\end{enumerate}
\end{ejemplo}

\begin{ejemplo}[Ejemplo tipo PAES 3: Herencia mendeliana]
{\sffamily\footnotesize\color{zinc500} Enunciado}\par
El albinismo se hereda como un rasgo autosómico recesivo. Si una pareja de padres no albinos tiene ambos el genotipo heterocigoto $Aa$, ¿cuál es la probabilidad de que un hijo sea albino?
\alternativa{A}{0 \%}
\alternativa{B}{25 \%}
\alternativa{C}{50 \%}
\alternativa{D}{75 \%}
\clave{B}
{\sffamily\footnotesize\color{zinc500} Análisis}\par Un rasgo autosómico recesivo se expresa solo en homocigotos $aa$; el cruce $Aa \times Aa$ da una probabilidad de $\frac{1}{4}$ de hijos albinos.
{\sffamily\footnotesize\color{zinc500} Desarrollo}
\begin{enumerate}
  \item Los padres son $Aa \times Aa$. Como el albinismo es recesivo, un individuo solo es albino si tiene genotipo $aa$.
  \item Cada padre produce gametos $A$ o $a$, con probabilidad $\frac{1}{2}$ cada uno, pues los alelos se segregan de forma independiente en la meiosis.
  \item En el cuadro de Punnett, las combinaciones posibles son $AA$, $Aa$, $Aa$ y $aa$, cada una con probabilidad $\frac{1}{4}$.
  \item Las combinaciones $AA$ y $Aa$ no son albinas, porque tienen al menos un alelo $A$. Solo la combinación $aa$ es albina.
  \item Entonces la probabilidad de albinismo es $P = \frac{1}{4} = 0{,}25$, es decir, $25\,\%$.
  \item Control: $\frac{1}{4} + \frac{2}{4} + \frac{1}{4} = 1$, así que las probabilidades cubren todas las combinaciones posibles.
  \item Así, un hijo tiene un 25 \% de probabilidad de ser albino. Por lo tanto, la alternativa correcta es la B.
\end{enumerate}
\end{ejemplo}

\subsection{Sistema nervioso}

El sistema nervioso recibe estímulos, los integra y coordina las respuestas del organismo. La neurona tiene dendritas (reciben), soma y axón (conduce el impulso); el impulso nervioso es un cambio eléctrico que viaja por el axón y pasa a la siguiente célula en la sinapsis química, mediante neurotransmisores.

En el arco reflejo simple, el estímulo va de un receptor a la médula espinal por una neurona sensitiva y vuelve a un músculo por una motora, sin pasar por el cerebro. El café, el alcohol y las drogas alteran la sinapsis; dormir lo suficiente y protegerse de golpes cuidan el sistema nervioso.

\begin{mecanismo}[Arco reflejo]
Receptor $\rightarrow$ neurona sensitiva (aferente) $\rightarrow$ médula espinal (interneurona) $\rightarrow$ neurona motora (eferente) $\rightarrow$ efector (músculo o glándula).
\end{mecanismo}

\begin{ideasclave}
\begin{itemize}
  \item En la sinapsis química el mensaje pasa con neurotransmisores.
  \item El reflejo es rápido porque no lo procesa el cerebro.
  \item Las drogas de abuso actúan sobre los neurotransmisores.
\end{itemize}
\end{ideasclave}

\begin{metodo}
\begin{enumerate}
  \item Ubica el estímulo, la vía y la respuesta.
  \item Identifica la estructura de la neurona o la etapa del arco reflejo involucrada.
  \item Predice qué pasa si esa estructura falla o una sustancia la altera.
\end{enumerate}
\end{metodo}

\begin{errorfrecuente}[Trampas frecuentes del DEMRE]
\begin{itemize}
  \item Pensar que todo reflejo pasa por el cerebro.
  \item Confundir dendrita (recibe) con axón (conduce el impulso).
\end{itemize}
\end{errorfrecuente}

\begin{ejemplo}[Ejemplo resuelto]
{\sffamily\footnotesize\color{zinc500} Enunciado}\par
Una persona toca una superficie caliente y retira la mano antes de sentir dolor. ¿Qué explica esa rapidez?
\alternativa{A}{El estímulo llega primero a la corteza cerebral.}
\alternativa{B}{La respuesta se integra en la médula espinal, en un arco reflejo.}
\alternativa{C}{Las hormonas actúan más rápido que los nervios.}
\alternativa{D}{La piel mueve los músculos por sí sola.}
\clave{B}
{\sffamily\footnotesize\color{zinc500} Análisis}\par Es un reflejo: la médula integra la respuesta sin esperar al cerebro.
{\sffamily\footnotesize\color{zinc500} Desarrollo}
\begin{enumerate}
  \item La neurona sensitiva lleva el estímulo a la médula.
  \item Desde ahí, una neurona motora activa el músculo; el dolor se percibe después.
\end{enumerate}
\end{ejemplo}

\subsection{Reproducción, sexualidad e ITS}

La fecundación es la unión de un ovocito y un espermatozoide, ambos haploides, que forma un cigoto diploide. El ciclo ovárico (fase folicular, ovulación y fase lútea) y el ciclo uterino (menstruación, proliferación y fase secretora) se coordinan cada unos 28 días.

Los métodos de control de la natalidad pueden ser naturales (Billings, calendario, temperatura basal), artificiales reversibles (hormonales y de barrera) o parcialmente reversibles (quirúrgicos). Solo los de barrera, como el condón, previenen además las infecciones de transmisión sexual (ITS) como VIH, herpes, gonorrea y clamidia. La sexualidad incluye aspectos biológicos, afectivos y sociales, con respeto y responsabilidad.

\begin{mecanismo}[Ciclo de 28 días]
Días 1–5: menstruación. Días 6–13: fase folicular, el endometrio se engrosa. Día 14 aprox.: ovulación. Días 15–28: fase lútea; si no hay fecundación, bajan las hormonas y comienza un nuevo ciclo.
\end{mecanismo}

\begin{ideasclave}
\begin{itemize}
  \item Ovulación: liberación del ovocito, cerca del día 14.
  \item Solo los métodos de barrera protegen de las ITS.
  \item El VIH es un virus; la gonorrea y la clamidia, bacterias.
\end{itemize}
\end{ideasclave}

\begin{metodo}
\begin{enumerate}
  \item Ubica en qué fase del ciclo ocurre lo descrito.
  \item Identifica el método o el agente infeccioso y cómo actúa.
  \item Decide si previene embarazo, ITS o ambos.
\end{enumerate}
\end{metodo}

\begin{errorfrecuente}[Trampas frecuentes del DEMRE]
\begin{itemize}
  \item Creer que los anticonceptivos hormonales protegen de las ITS.
  \item Ubicar la ovulación al inicio del ciclo.
\end{itemize}
\end{errorfrecuente}

\begin{ejemplo}[Ejemplo resuelto]
{\sffamily\footnotesize\color{zinc500} Enunciado}\par
¿Cuál de los siguientes métodos previene a la vez el embarazo y las infecciones de transmisión sexual?
\alternativa{A}{Píldora anticonceptiva.}
\alternativa{B}{Método del calendario.}
\alternativa{C}{Preservativo.}
\alternativa{D}{Ligadura de trompas.}
\clave{C}
{\sffamily\footnotesize\color{zinc500} Análisis}\par Solo una barrera física impide el contacto con fluidos y patógenos.
{\sffamily\footnotesize\color{zinc500} Desarrollo}
\begin{enumerate}
  \item La píldora y la ligadura evitan el embarazo, pero no las ITS.
  \item El calendario no protege de ninguna de las dos de forma confiable.
\end{enumerate}
\end{ejemplo}

\subsection{Ciclo celular, meiosis y manipulación genética}

El ciclo celular tiene una interfase ($G_1$, $S$ en que se replica el ADN, y $G_2$) y la mitosis (profase, metafase, anafase y telofase). Los puntos de control $G_1$–$S$, $G_2$–$M$ y de metafase revisan que todo esté en orden; si fallan, la célula puede dividirse sin control, como en el cáncer.

La mitosis produce dos células hijas idénticas a la madre (crecimiento, desarrollo y reparación de tejidos). La meiosis produce cuatro células haploides genéticamente distintas (gametos): el entrecruzamiento y la permutación cromosómica generan variabilidad.

La manipulación genética inserta, modifica o elimina genes de un organismo. Se usa para producir alimentos transgénicos, fármacos como la insulina humana en bacterias, detergentes con enzimas o fibras para vestuario.

\begin{mecanismo}[Mitosis y meiosis]
Mitosis: $2n\rightarrow$ 2 células $2n$ idénticas. Meiosis: $2n\rightarrow$ 4 células $n$ distintas.

Humanos: $2n=46$ en células somáticas y $n=23$ en gametos. La cromatina se compacta en cromosomas visibles durante la división.
\end{mecanismo}

{\sffamily\footnotesize\color{zinc500} Esquema comparativo}
\begin{center}\small
\begin{tabular}{@{}>{\raggedright\arraybackslash}p{\dimexpr(\linewidth-4\tabcolsep)/3\relax}>{\raggedright\arraybackslash}p{\dimexpr(\linewidth-4\tabcolsep)/3\relax}>{\raggedright\arraybackslash}p{\dimexpr(\linewidth-4\tabcolsep)/3\relax}@{}}
\toprule
 & \textbf{Mitosis} & \textbf{Meiosis} \\
\midrule
Células hijas & 2 & 4 \\
Cromosomas & igual a la madre ($2n$) & la mitad ($n$) \\
Función & crecimiento y reparación & formar gametos \\
\bottomrule
\end{tabular}
\end{center}

\begin{ideasclave}
\begin{itemize}
  \item La replicación del ADN ocurre en la fase $S$, no en la mitosis.
  \item La variabilidad viene de la meiosis, no de la mitosis.
  \item Un transgénico porta un gen de otra especie.
\end{itemize}
\end{ideasclave}

\begin{metodo}
\begin{enumerate}
  \item Identifica en qué fase del ciclo o de la división ocurre lo descrito.
  \item Cuenta cromosomas y moléculas de ADN en cada etapa.
  \item Relaciona el proceso con su función: crecimiento, reparación, gametos o aplicación biotecnológica.
\end{enumerate}
\end{metodo}

\begin{errorfrecuente}[Trampas frecuentes del DEMRE]
\begin{itemize}
  \item Atribuir a la mitosis la variabilidad genética.
  \item Confundir la duplicación del ADN (fase $S$) con la separación de cromátidas (anafase).
\end{itemize}
\end{errorfrecuente}

\begin{ejemplo}[Ejemplo resuelto]
{\sffamily\footnotesize\color{zinc500} Enunciado}\par
Una célula humana somática con 46 cromosomas completa una mitosis. ¿Cuántos cromosomas tiene cada célula hija?
\alternativa{A}{23}
\alternativa{B}{46}
\alternativa{C}{69}
\alternativa{D}{92}
\clave{B}
{\sffamily\footnotesize\color{zinc500} Análisis}\par La mitosis conserva el número de cromosomas.
{\sffamily\footnotesize\color{zinc500} Desarrollo}
\begin{enumerate}
  \item Antes de dividirse, cada cromosoma se duplicó en la fase $S$.
  \item En la anafase se separan las cromátidas y cada hija recibe 46 cromosomas.
\end{enumerate}
\end{ejemplo}

\begin{ejemplo}[Ejemplo tipo PAES 2: División celular]
{\sffamily\footnotesize\color{zinc500} Enunciado}\par
Una célula humana con 46 cromosomas ($2n = 46$) completa una meiosis. ¿Qué se obtiene al final del proceso?
\alternativa{A}{4 células con 23 cromosomas cada una.}
\alternativa{B}{2 células con 46 cromosomas cada una.}
\alternativa{C}{4 células con 46 cromosomas cada una.}
\alternativa{D}{2 células con 23 cromosomas cada una.}
\clave{A}
{\sffamily\footnotesize\color{zinc500} Análisis}\par En la meiosis, dos divisiones sucesivas reducen a la mitad el número de cromosomas: de una célula diploide se obtienen cuatro células haploides.
{\sffamily\footnotesize\color{zinc500} Desarrollo}
\begin{enumerate}
  \item La meiosis I separa los cromosomas homólogos. Al final de esta división cada célula hija tiene un juego de cromosomas, es decir, $n = 23$, aunque cada cromosoma aún conserva sus dos cromátidas.
  \item Así, se obtienen dos células con $n = 23$ cromosomas cada una, y ambas siguen en proceso de división.
  \item La meiosis II separa las cromátidas hermanas, de manera parecida a una mitosis, pero sobre células haploides. Cada una de las dos células se divide en dos.
  \item Resultado: $2 \times 2 = 4$ células, cada una con $n = 23$ cromosomas simples.
  \item Control: la célula original replicó su ADN antes de la meiosis, por lo que el total de cromosomas simples en las cuatro células es $4 \cdot 23 = 92$, el doble de los 46 de la célula inicial.
  \item La alternativa B (2 células con 46 cromosomas) corresponde al resultado de una mitosis, que conserva el número cromosómico y produce dos células.
  \item La alternativa C mantiene el número de cromosomas de la célula madre, lo que no ocurre en la meiosis, que reduce a la mitad.
  \item La alternativa D describe el estado al final de la meiosis I, no el producto final del proceso completo.
  \item Así, al completar la meiosis se obtienen cuatro células haploides con 23 cromosomas cada una. Por lo tanto, la alternativa correcta es la A.
\end{enumerate}
\end{ejemplo}

\subsection{Evolución: evidencias y selección natural}

La teoría de la evolución por selección natural (Darwin y Wallace) explica que en una población hay variabilidad heredable; los individuos con características ventajosas en un ambiente dejan más descendencia y esas características se vuelven más frecuentes.

Las evidencias de la evolución incluyen el registro fósil, las estructuras homólogas, la embriología comparada y las semejanzas en el ADN.

\begin{mecanismo}[Selección natural]
Variabilidad $\rightarrow$ presión del ambiente $\rightarrow$ supervivencia y reproducción diferencial $\rightarrow$ cambio en la frecuencia de los rasgos en la población.
\end{mecanismo}

\begin{ideasclave}
\begin{itemize}
  \item Evolucionan las poblaciones, no los individuos.
  \item Las mutaciones ocurren al azar; la selección no es azarosa.
  \item Lamarck (caracteres adquiridos) fue descartado.
\end{itemize}
\end{ideasclave}

\begin{metodo}
\begin{enumerate}
  \item Identifica la variación presente en la población.
  \item Determina qué presión ambiental favorece a qué variante.
  \item Concluye cómo cambia la frecuencia del rasgo en generaciones.
\end{enumerate}
\end{metodo}

\begin{errorfrecuente}[Trampas frecuentes del DEMRE]
\begin{itemize}
  \item Pensar que el individuo cambia «para adaptarse».
  \item Explicar con herencia de caracteres adquiridos.
\end{itemize}
\end{errorfrecuente}

\begin{ejemplo}[Ejemplo resuelto]
{\sffamily\footnotesize\color{zinc500} Enunciado}\par
Tras usar un antibiótico por años, en un hospital predominan bacterias resistentes. ¿Cuál es la mejor explicación?
\alternativa{A}{El antibiótico produjo mutaciones dirigidas a la resistencia.}
\alternativa{B}{Las bacterias resistentes ya existían y sobrevivieron y se reprodujeron más.}
\alternativa{C}{Las bacterias aprendieron a resistir el antibiótico.}
\alternativa{D}{El antibiótico perdió su efecto por el paso del tiempo.}
\clave{B}
{\sffamily\footnotesize\color{zinc500} Análisis}\par Es un ejemplo clásico de selección natural.
{\sffamily\footnotesize\color{zinc500} Desarrollo}
\begin{enumerate}
  \item La resistencia surgió por variación previa (mutaciones al azar).
  \item El antibiótico eliminó a las sensibles y las resistentes dejaron más descendencia.
\end{enumerate}
\end{ejemplo}

\subsection{Fotosíntesis, respiración y flujo de energía}

Los organismos autótrofos fabrican materia orgánica: en la fotosíntesis, que ocurre en los cloroplastos, la luz, el $\mathrm{CO_2}$ y el agua producen glucosa y oxígeno. La etapa dependiente de la luz ocurre en los tilacoides (libera $\mathrm{O_2}$ y forma ATP); la independiente, en el estroma (fija el $\mathrm{CO_2}$ en el ciclo de Calvin).

La respiración celular, en las mitocondrias, degrada la glucosa con oxígeno y libera energía (ATP), $\mathrm{CO_2}$ y agua. Así se reciclan el carbono y el oxígeno, mientras la energía fluye en una sola dirección y en cada nivel trófico pasa solo cerca del 10 \% al siguiente.

La intensidad de la luz, la concentración de $\mathrm{CO_2}$ y la temperatura limitan la rapidez de la fotosíntesis.

\begin{mecanismo}[Ecuaciones globales]
\[6\,\mathrm{CO_2}+6\,\mathrm{H_2O}\xrightarrow{\text{luz}}\mathrm{C_6H_{12}O_6}+6\,\mathrm{O_2}\]

Respiración: la reacción inversa, que libera la energía de la glucosa como ATP.
\end{mecanismo}

\begin{ideasclave}
\begin{itemize}
  \item Las plantas también respiran, día y noche.
  \item El $\mathrm{O_2}$ liberado proviene del agua.
  \item La energía fluye; la materia se recicla.
\end{itemize}
\end{ideasclave}

\begin{metodo}
\begin{enumerate}
  \item Identifica si el proceso fabrica materia orgánica o la degrada.
  \item Ubica la etapa y el organelo (cloroplasto o mitocondria).
  \item Predice el efecto de cambiar la luz, el $\mathrm{CO_2}$ o la temperatura.
\end{enumerate}
\end{metodo}

\begin{errorfrecuente}[Trampas frecuentes del DEMRE]
\begin{itemize}
  \item Creer que las plantas no respiran.
  \item Decir que la energía se recicla en el ecosistema.
\end{itemize}
\end{errorfrecuente}

\begin{ejemplo}[Ejemplo resuelto]
{\sffamily\footnotesize\color{zinc500} Enunciado}\par
En un experimento, una planta acuática se ilumina con luz cada vez más intensa. ¿Qué es lo más probable que ocurra con la liberación de burbujas de oxígeno?
\alternativa{A}{Disminuye siempre.}
\alternativa{B}{Aumenta hasta un punto y luego se estabiliza.}
\alternativa{C}{No cambia.}
\alternativa{D}{Aumenta sin límite.}
\clave{B}
{\sffamily\footnotesize\color{zinc500} Análisis}\par La luz es un factor limitante solo hasta que otro factor limita el proceso.
{\sffamily\footnotesize\color{zinc500} Desarrollo}
\begin{enumerate}
  \item Más luz acelera la etapa luminosa y la liberación de $\mathrm{O_2}$.
  \item Luego el $\mathrm{CO_2}$ o la temperatura limitan y la tasa se estabiliza.
\end{enumerate}
\end{ejemplo}

\begin{ejemplo}[Ejemplo tipo PAES 2: Fotosíntesis]
{\sffamily\footnotesize\color{zinc500} Enunciado}\par
Se riega una planta con agua cuyo oxígeno es el isótopo $^{18}\mathrm{O}$ y se la ilumina para que realice fotosíntesis. ¿Dónde se detectará principalmente el $^{18}\mathrm{O}$?
\alternativa{A}{En la glucosa producida.}
\alternativa{B}{En el $\mathrm{CO_2}$ absorbido por las hojas.}
\alternativa{C}{En la clorofila de los cloroplastos.}
\alternativa{D}{En el $\mathrm{O_2}$ liberado a la atmósfera.}
\clave{D}
{\sffamily\footnotesize\color{zinc500} Análisis}\par En la fase luminosa de la fotosíntesis, la fotólisis del agua libera su oxígeno como $\mathrm{O_2}$; por eso el isótopo marcado aparece en el oxígeno liberado.
{\sffamily\footnotesize\color{zinc500} Desarrollo}
\begin{enumerate}
  \item La ecuación global de la fotosíntesis es $6\,\mathrm{CO_2} + 6\,\mathrm{H_2O} \xrightarrow{\text{luz}} \mathrm{C_6H_{12}O_6} + 6\,\mathrm{O_2}$. El oxígeno que se libera como $\mathrm{O_2}$ proviene del agua.
  \item En la fase dependiente de la luz el agua se fotoliza, es decir, se rompe: $2\,\mathrm{H_2O} \rightarrow 4\,\mathrm{H^+} + 4\,e^- + \mathrm{O_2}$. El oxígeno del agua pasa al $\mathrm{O_2}$ liberado.
  \item Por lo tanto, los átomos de oxígeno que quedan como $\mathrm{O_2}$ provienen del agua, y el isótopo marcado sale a la atmósfera en esa forma.
  \item El $\mathrm{CO_2}$ que la planta fija no lleva el isótopo marcado, y los átomos de oxígeno de la glucosa provienen principalmente del $\mathrm{CO_2}$ fijado en el ciclo de Calvin. Por eso el marcado no se concentra en la glucosa.
  \item La clorofila es el pigmento que capta la luz; no es un producto de la reacción ni almacena el oxígeno del agua, así que no es el destino del isótopo.
  \item Así, el $^{18}\mathrm{O}$ se detecta principalmente en el $\mathrm{O_2}$ liberado a la atmósfera. Por lo tanto, la alternativa correcta es la D.
\end{enumerate}
\end{ejemplo}

\begin{ejemplo}[Ejemplo tipo PAES 3: Flujo de energía]
{\sffamily\footnotesize\color{zinc500} Enunciado}\par
En un ecosistema, los productores fijan 20 000 kJ de energía. Si en cada nivel trófico se transfiere aproximadamente el 10 \% de la energía al nivel siguiente, ¿cuánta energía llega a los consumidores terciarios?
\alternativa{A}{2 kJ}
\alternativa{B}{20 kJ}
\alternativa{C}{200 kJ}
\alternativa{D}{2000 kJ}
\clave{B}
{\sffamily\footnotesize\color{zinc500} Análisis}\par En cada transferencia entre niveles tróficos se conserva cerca del 10 \% de la energía; por eso los consumidores terciarios reciben $20$ kJ de los $20\,000$ kJ fijados.
{\sffamily\footnotesize\color{zinc500} Desarrollo}
\begin{enumerate}
  \item Los productores fijan $20\,000$ kJ; ese es el punto de partida de la cadena.
  \item Consumidores primarios (herbívoros): $20\,000 \cdot 0{,}1 = 2000$ kJ.
  \item Consumidores secundarios (carnívoros de primer orden): $2000 \cdot 0{,}1 = 200$ kJ.
  \item Consumidores terciarios (carnívoros de segundo orden): $200 \cdot 0{,}1 = 20$ kJ.
  \item Control: cada paso multiplica la energía por $0{,}1$, así que tres transferencias dan $20\,000 \cdot 0{,}1^3 = 20$ kJ.
  \item La alternativa D (2000 kJ) es la energía de los consumidores primarios; la alternativa C (200 kJ) corresponde a los secundarios.
  \item La alternativa A (2 kJ) requiere una transferencia más, y corresponde a los consumidores cuaternarios.
  \item Así, los consumidores terciarios reciben 20 kJ. Por lo tanto, la alternativa correcta es la B.
\end{enumerate}
\end{ejemplo}

\section{Física}

{\itshape\color{zinc700} Ondas electromagnéticas, mecánica (movimiento, fuerzas y universo), energía y Tierra, y electricidad.}

\subsection{Ondas electromagnéticas, luz y óptica}

Las ondas electromagnéticas se propagan también en el vacío, a $c=3\cdot 10^8$ m/s. Se describen con su longitud de onda $\lambda$, frecuencia $f$, período $T$ y amplitud, y cumplen $c=\lambda f$. Ordenadas por frecuencia forman el espectro: ondas de radio, microondas, infrarrojo, luz visible, ultravioleta, rayos X y gamma.

La luz se propaga en línea recta y puede absorberse, reflejarse o refractarse. La refracción en un prisma separa los colores (dispersión); el color de un objeto es la luz que refleja. La interferencia, la difracción y el efecto Doppler (cambio de frecuencia por el movimiento relativo) se estudian de forma cualitativa.

Los espejos planos forman imágenes virtuales del mismo tamaño; los cóncavos pueden formar imágenes reales o virtuales y los convexos, siempre virtuales y menores. Las lentes convergentes concentran la luz y las divergentes la dispersan. Así funcionan telescopios, prismáticos, fibra óptica, radares y teléfonos.

\begin{ley}[Relaciones]
\begin{gather*}
c=\lambda\,f \\
T=\frac{1}{f} \\
\theta_i=\theta_r \\
n=\frac{c}{v}
\end{gather*}

Más frecuencia $\Leftrightarrow$ menor longitud de onda $\Leftrightarrow$ más energía (radio $<$ microondas $<$ infrarrojo $<$ visible $<$ UV $<$ X $<$ gamma).
\end{ley}

\begin{ideasclave}
\begin{itemize}
  \item Todas las ondas electromagnéticas viajan a la misma rapidez en el vacío.
  \item Los ángulos de reflexión y refracción se miden respecto de la normal.
  \item El arcoíris se explica por refracción y dispersión.
\end{itemize}
\end{ideasclave}

\begin{metodo}
\begin{enumerate}
  \item Identifica qué magnitud de la onda se pregunta y cuáles conoces.
  \item Ubica la radiación en el espectro o dibuja la normal y los rayos.
  \item Aplica $c=\lambda f$, la ley de reflexión o el tipo de imagen según el espejo o la lente.
\end{enumerate}
\end{metodo}

\begin{errorfrecuente}[Trampas frecuentes del DEMRE]
\begin{itemize}
  \item Pensar que la frecuencia cambia al pasar de un medio a otro (cambian la rapidez y la longitud de onda).
  \item Creer que los rayos X o la luz necesitan un medio para propagarse.
\end{itemize}
\end{errorfrecuente}

\begin{ejemplo}[Ejemplo resuelto]
{\sffamily\footnotesize\color{zinc500} Enunciado}\par
Una radio transmite con una frecuencia de $100\ \text{MHz}=1\cdot 10^8$ Hz. ¿Cuál es la longitud de onda de esa señal en el aire?
\alternativa{A}{0,3 m}
\alternativa{B}{1 m}
\alternativa{C}{3 m}
\alternativa{D}{30 m}
\clave{C}
{\sffamily\footnotesize\color{zinc500} Análisis}\par Las ondas de radio son electromagnéticas: viajan a unos $3\cdot 10^8$ m/s.
{\sffamily\footnotesize\color{zinc500} Desarrollo}
\begin{enumerate}
  \item $\lambda=\dfrac{c}{f}=\dfrac{3\cdot 10^8}{1\cdot 10^8}$.
  \item $\lambda=3$ m.
\end{enumerate}
\end{ejemplo}

\begin{ejemplo}[Ejemplo tipo PAES 2: Ondas]
{\sffamily\footnotesize\color{zinc500} Enunciado}\par
Un sonido de 680 Hz se propaga en el aire con una rapidez de 340 m/s. ¿Cuál es su longitud de onda?
\alternativa{A}{0,5 m}
\alternativa{B}{2 m}
\alternativa{C}{1020 m}
\alternativa{D}{231 200 m}
\clave{A}
{\sffamily\footnotesize\color{zinc500} Análisis}\par Para cualquier onda se cumple $v = \lambda f$; despejando, la longitud de onda es el cociente entre la rapidez y la frecuencia.
{\sffamily\footnotesize\color{zinc500} Desarrollo}
\begin{enumerate}
  \item La relación fundamental de las ondas es $v = \lambda\, f$, donde $\lambda$ es la longitud de onda y $f$ la frecuencia.
  \item Se despeja $\lambda$ dividiendo ambos lados por $f$: $\lambda = \dfrac{v}{f}$.
  \item Se reemplazan los valores: $\lambda = \dfrac{340\ \mathrm{m/s}}{680\ \mathrm{s^{-1}}}$. Las unidades $\mathrm{m/s}$ divididas entre $\mathrm{s^{-1}}$ dan metros.
  \item Simplificando el cociente, $\frac{340}{680} = \frac{1}{2}$, así que $\lambda = 0{,}5\ \mathrm{m}$.
  \item Control: con $\lambda = 0{,}5$ m y $f = 680$ Hz, la rapidez es $0{,}5 \cdot 680 = 340$ m/s, que es el valor del enunciado.
  \item La alternativa B (2 m) invierte el cociente: $f/v = 680/340 = 2$, pero con unidades de $\mathrm{m^{-1}}$, así que no es una longitud.
  \item La alternativa C (1020 m) suma $v + f$, y la D (231 200 m) multiplica $v \cdot f$; ninguna de esas operaciones corresponde a la relación de ondas.
  \item Así, la longitud de onda es de $0{,}5$ m. Por lo tanto, la alternativa correcta es la A.
\end{enumerate}
\end{ejemplo}

\begin{ejemplo}[Ejemplo tipo PAES 3: Conservación de la energía]
{\sffamily\footnotesize\color{zinc500} Enunciado}\par
Una pelota de 2 kg se deja caer desde 5 m de altura. Si se desprecia el roce con el aire y $g = 10\ \mathrm{m/s^2}$, ¿con qué rapidez llega al suelo?
\alternativa{A}{5 m/s}
\alternativa{B}{50 m/s}
\alternativa{C}{100 m/s}
\alternativa{D}{10 m/s}
\clave{D}
{\sffamily\footnotesize\color{zinc500} Análisis}\par Sin roce, la energía mecánica se conserva: la energía potencial inicial se transforma en energía cinética, y la masa se cancela al despejar la rapidez.
{\sffamily\footnotesize\color{zinc500} Desarrollo}
\begin{enumerate}
  \item En la altura inicial la pelota tiene solo energía potencial gravitatoria, $E_p = mgh$, y parte del reposo, por lo que su energía cinética es cero.
  \item Sin roce con el aire, la energía mecánica se conserva: la energía potencial inicial se transforma en energía cinética al llegar al suelo.
  \item Se calcula la energía inicial: $E_p = 2\ \mathrm{kg}\cdot 10\ \mathrm{m/s^2}\cdot 5\ \mathrm{m} = 100\ \mathrm{J}$.
  \item Se iguala con la energía cinética final: $\tfrac{1}{2}mv^2 = 100\ \mathrm{J}$, es decir, $\tfrac{1}{2}\cdot 2\cdot v^2 = 100$.
  \item Simplificando, $v^2 = 100\ \mathrm{m^2/s^2}$, y al tomar la raíz positiva resulta $v = 10\ \mathrm{m/s}$.
  \item Control: la fórmula general $v = \sqrt{2gh} = \sqrt{2\cdot 10\cdot 5} = \sqrt{100} = 10$ m/s coincide. La masa no aparece porque se cancela en ambos lados.
  \item La alternativa A (5 m/s) toma la altura en metros como si fuera una rapidez; la alternativa B (50 m/s) no proviene de ninguna relación correcta, como $gh$ o $\sqrt{2gh}$.
  \item La alternativa C (100 m/s) corresponde a la energía en joules o al valor de $v^2$, y no a la rapidez.
  \item Así, la pelota llega al suelo con $10$ m/s. Por lo tanto, la alternativa correcta es la D.
\end{enumerate}
\end{ejemplo}

\subsection{Movimiento: MRU y MRUA}

La cinemática describe el movimiento. La rapidez media es la distancia recorrida dividida por el tiempo; la velocidad considera además la dirección. La aceleración mide cuánto cambia la velocidad por unidad de tiempo.

En el movimiento rectilíneo uniforme (MRU) la velocidad es constante; en el uniformemente acelerado (MRUA), la aceleración es constante. En un gráfico posición–tiempo la pendiente es la velocidad; en uno velocidad–tiempo, la pendiente es la aceleración y el área bajo la curva, el desplazamiento.

\begin{formula}[Ecuaciones]
\begin{gather*}
\text{MRU: } x=x_0+vt \\
\text{MRUA: } v=v_0+at,\quad x=x_0+v_0t+\tfrac{1}{2}at^2
\end{gather*}
\end{formula}

\begin{ideasclave}
\begin{itemize}
  \item Pendiente de $x$–$t$ = velocidad.
  \item Área bajo $v$–$t$ = desplazamiento.
  \item 72 km/h = 20 m/s (divide por 3,6).
\end{itemize}
\end{ideasclave}

\begin{metodo}
\begin{enumerate}
  \item Identifica el tipo de movimiento y los datos con sus unidades.
  \item Elige la ecuación que relaciona datos e incógnita.
  \item En gráficos, usa la pendiente o el área según lo pedido.
\end{enumerate}
\end{metodo}

\begin{errorfrecuente}[Trampas frecuentes del DEMRE]
\begin{itemize}
  \item Mezclar km/h con m/s.
  \item Olvidar el factor $\frac{1}{2}$ en $\frac{1}{2}at^2$.
\end{itemize}
\end{errorfrecuente}

\begin{ejemplo}[Ejemplo resuelto]
{\sffamily\footnotesize\color{zinc500} Enunciado}\par
Un ciclista recorre 18 km en 45 minutos. ¿Cuál es su rapidez media?
\alternativa{A}{0,4 km/h}
\alternativa{B}{2,5 km/h}
\alternativa{C}{24 km/h}
\alternativa{D}{40 km/h}
\clave{C}
{\sffamily\footnotesize\color{zinc500} Análisis}\par Rapidez media = distancia / tiempo, con unidades coherentes.
{\sffamily\footnotesize\color{zinc500} Desarrollo}
\begin{enumerate}
  \item 45 min = 0,75 h.
  \item $v=\dfrac{18}{0{,}75}=24$ km/h.
\end{enumerate}
\end{ejemplo}

\begin{ejemplo}[Ejemplo tipo PAES 2: Movimiento rectilíneo uniformemente acelerado]
{\sffamily\footnotesize\color{zinc500} Enunciado}\par
Un automóvil parte del reposo y acelera de manera constante a $2\ \mathrm{m/s^2}$ durante 5 s. ¿Qué distancia recorre en ese tiempo?
\alternativa{A}{5 m}
\alternativa{B}{10 m}
\alternativa{C}{25 m}
\alternativa{D}{50 m}
\clave{C}
{\sffamily\footnotesize\color{zinc500} Análisis}\par En un movimiento uniformemente acelerado desde el reposo, la distancia es $d = \tfrac{1}{2}at^2$; con $a = 2\ \mathrm{m/s^2}$ y $t = 5$ s resulta $25$ m.
{\sffamily\footnotesize\color{zinc500} Desarrollo}
\begin{enumerate}
  \item Como el auto parte del reposo, $v_0 = 0$, y la aceleración es constante. La distancia recorrida se describe con $d = v_0 t + \tfrac{1}{2}at^2$.
  \item Al anular el primer término por ser $v_0 = 0$, queda $d = \tfrac{1}{2}at^2$.
  \item Reemplazando: $d = \tfrac{1}{2}\cdot 2\ \mathrm{m/s^2}\cdot (5\ \mathrm{s})^2 = 1\ \mathrm{m/s^2}\cdot 25\ \mathrm{s^2} = 25\ \mathrm{m}$.
  \item Control: la rapidez final es $v = at = 2 \cdot 5 = 10$ m/s. Como la aceleración es constante, la rapidez media es $\frac{0 + 10}{2} = 5$ m/s, y $d = \bar{v}\,t = 5\cdot 5 = 25$ m.
  \item La alternativa B (10 m) es la rapidez final, $v = at$, y no la distancia.
  \item La alternativa A (5 m) corresponde a $\tfrac{1}{2}at$, es decir, usar el tiempo en lugar de su cuadrado.
  \item La alternativa D (50 m) se obtiene al omitir el factor $\tfrac{1}{2}$ de la fórmula.
  \item Así, el automóvil recorre 25 m. Por lo tanto, la alternativa correcta es la C.
\end{enumerate}
\end{ejemplo}

\subsection{Fuerzas y leyes de Newton}

Primera ley (inercia): si la fuerza neta es cero, un cuerpo mantiene su estado de reposo o de velocidad constante. Segunda ley: $F_{\text{neta}}=ma$. Tercera ley: a toda acción corresponde una reacción de igual magnitud y sentido opuesto, aplicada sobre el \textbf{otro} cuerpo.

El peso es la fuerza de gravedad ($P=mg$) y depende del lugar; la masa no. El roce estático impide que un cuerpo empiece a deslizar y el cinético se opone al deslizamiento; la fuerza elástica de un resorte cumple la ley de Hooke, $F=kx$.

\begin{ley}[Fórmulas]
\begin{gather*}
F_{\text{neta}}=m\,a \\
P=m\,g \\
F_{\text{elástica}}=k\,x \\
f_r=\mu\,N
\end{gather*}
\end{ley}

\begin{ideasclave}
\begin{itemize}
  \item Acción y reacción actúan sobre cuerpos distintos: no se anulan.
  \item Velocidad constante $\Rightarrow$ fuerza neta nula.
  \item En la Luna cambia el peso, no la masa.
\end{itemize}
\end{ideasclave}

\begin{metodo}
\begin{enumerate}
  \item Dibuja las fuerzas que actúan sobre el cuerpo (diagrama de cuerpo libre).
  \item Calcula la fuerza neta sumando con signo.
  \item Aplica $F=ma$ o la ley pertinente.
\end{enumerate}
\end{metodo}

\begin{errorfrecuente}[Trampas frecuentes del DEMRE]
\begin{itemize}
  \item Creer que se necesita una fuerza para mantener la velocidad constante.
  \item Sumar acción y reacción como si actuaran sobre el mismo cuerpo.
\end{itemize}
\end{errorfrecuente}

\begin{ejemplo}[Ejemplo resuelto]
{\sffamily\footnotesize\color{zinc500} Enunciado}\par
Un auto avanza en línea recta a 60 km/h constantes. ¿Qué se puede afirmar de la fuerza neta sobre él?
\alternativa{A}{Es nula.}
\alternativa{B}{Apunta en el sentido del movimiento.}
\alternativa{C}{Apunta en sentido contrario al movimiento.}
\alternativa{D}{Es igual a su peso.}
\clave{A}
{\sffamily\footnotesize\color{zinc500} Análisis}\par Velocidad constante significa aceleración cero.
{\sffamily\footnotesize\color{zinc500} Desarrollo}
\begin{enumerate}
  \item Si $a=0$, por la segunda ley $F_{\text{neta}}=m\cdot 0=0$.
  \item La fuerza del motor se equilibra con los roces.
\end{enumerate}
\end{ejemplo}

\begin{ejemplo}[Ejemplo tipo PAES 2: Segunda ley de Newton]
{\sffamily\footnotesize\color{zinc500} Enunciado}\par
Sobre una caja de 4 kg que se desliza por el suelo actúa una fuerza horizontal de 20 N. La fuerza de roce cinético con el suelo es de 4 N. ¿Cuál es la magnitud de la aceleración de la caja?
\alternativa{A}{$4\ \mathrm{m/s^2}$}
\alternativa{B}{$5\ \mathrm{m/s^2}$}
\alternativa{C}{$6\ \mathrm{m/s^2}$}
\alternativa{D}{$16\ \mathrm{m/s^2}$}
\clave{A}
{\sffamily\footnotesize\color{zinc500} Análisis}\par La aceleración depende de la fuerza neta: el roce se opone al movimiento y resta a los 20 N, dejando 16 N para acelerar la masa de 4 kg.
{\sffamily\footnotesize\color{zinc500} Desarrollo}
\begin{enumerate}
  \item Se dibuja el diagrama de fuerzas horizontales: la fuerza aplicada apunta en el sentido del movimiento y el roce cinético apunta en sentido contrario.
  \item La fuerza neta horizontal es la suma vectorial: $F_{\text{neta}} = F - f_k = 20\ \mathrm{N} - 4\ \mathrm{N} = 16\ \mathrm{N}$.
  \item En la vertical, el peso y la normal se equilibran, así que no contribuyen a la aceleración horizontal.
  \item Se aplica la segunda ley de Newton, $F_{\text{neta}} = m\,a$, de donde $a = \dfrac{F_{\text{neta}}}{m}$.
  \item Reemplazando: $a = \dfrac{16\ \mathrm{N}}{4\ \mathrm{kg}} = 4\ \mathrm{m/s^2}$.
  \item Control: con $a = 4\ \mathrm{m/s^2}$, la fuerza neta es $4\cdot 4 = 16$ N, que coincide con la diferencia entre 20 N y 4 N.
  \item La alternativa B (5 m/s\textsuperscript{2}) ignora el roce, pues $20/4 = 5$; la C (6 m/s\textsuperscript{2}) suma el roce en lugar de restarlo, pues $(20 + 4)/4 = 6$.
  \item La alternativa D (16 m/s\textsuperscript{2}) es la fuerza neta en newtons, no la aceleración.
  \item Así, la aceleración de la caja es de $4\ \mathrm{m/s^2}$. Por lo tanto, la alternativa correcta es la A.
\end{enumerate}
\end{ejemplo}

\subsection{Sistema solar y universo}

El modelo geocéntrico de Aristóteles y Ptolomeo ponía a la Tierra en el centro; el heliocéntrico de Copérnico, al Sol. Galileo lo apoyó con el telescopio (las lunas de Júpiter y las fases de Venus) y Kepler describió las órbitas como elipses, con el Sol en uno de sus focos.

La teoría del Big Bang propone que el universo se originó hace unos 13 800 millones de años desde un estado muy denso y caliente y que sigue expandiéndose. La sustentan el alejamiento de las galaxias (corrimiento al rojo) y la radiación de fondo de microondas. El Big Crunch plantea un posible colapso futuro.

\begin{ley}[Leyes de Kepler]
1.ª: los planetas giran en órbitas elípticas con el Sol en un foco. 2.ª: la línea Sol–planeta barre áreas iguales en tiempos iguales (el planeta va más rápido cerca del Sol). 3.ª: el cuadrado del período es proporcional al cubo del semieje mayor, $T^2\propto a^3$.
\end{ley}

\begin{ideasclave}
\begin{itemize}
  \item Las lunas de Júpiter mostraron que no todo gira en torno a la Tierra.
  \item Un planeta es más rápido en su perihelio que en su afelio.
  \item Corrimiento al rojo: la galaxia se aleja.
\end{itemize}
\end{ideasclave}

\begin{metodo}
\begin{enumerate}
  \item Identifica qué modelo o teoría describe la situación.
  \item Busca la evidencia que la apoya o la refuta.
  \item Relaciona la observación con la ley o el científico correspondiente.
\end{enumerate}
\end{metodo}

\begin{errorfrecuente}[Trampas frecuentes del DEMRE]
\begin{itemize}
  \item Atribuir a Copérnico las órbitas elípticas, que son de Kepler.
  \item Pensar que el Big Bang fue una explosión en un punto del espacio y no una expansión del espacio mismo.
\end{itemize}
\end{errorfrecuente}

\begin{ejemplo}[Ejemplo resuelto]
{\sffamily\footnotesize\color{zinc500} Enunciado}\par
¿Qué observación de Galileo con su telescopio contradijo el modelo geocéntrico?
\alternativa{A}{Que la Tierra es redonda.}
\alternativa{B}{Que Júpiter tiene lunas que giran a su alrededor.}
\alternativa{C}{Que las órbitas son elípticas.}
\alternativa{D}{Que el universo se expande.}
\clave{B}
{\sffamily\footnotesize\color{zinc500} Análisis}\par El modelo geocéntrico afirmaba que todos los astros giraban en torno a la Tierra.
{\sffamily\footnotesize\color{zinc500} Desarrollo}
\begin{enumerate}
  \item Las lunas de Júpiter giran en torno a otro cuerpo, no a la Tierra.
  \item Las órbitas elípticas son de Kepler, y la expansión del universo es del siglo XX.
\end{enumerate}
\end{ejemplo}

\subsection{Tectónica de placas y sismos}

La geosfera tiene corteza, manto y núcleo. La litosfera (corteza y parte superior del manto) está dividida en placas que se desplazan sobre la astenosfera, impulsadas por corrientes de convección del manto.

La deriva continental de Wegener se apoyó en el encaje de los continentes, fósiles y rocas iguales a ambos lados del océano. La tectónica de placas explica que en los bordes convergentes (como la subducción de Nazca bajo la Sudamericana) se generen sismos, volcanes y cordilleras; en los divergentes, nuevo suelo oceánico.

En un sismo, el hipocentro es el punto interior donde se libera la energía y el epicentro, el punto de la superficie justo sobre él. La magnitud mide la energía liberada; la intensidad, los efectos en cada lugar.

\begin{definicion}[Bordes de placa]
Convergente: las placas chocan; subducción, sismos, volcanes y cordilleras (Chile).

Divergente: se separan; dorsales y nuevo suelo oceánico. Transformante: se deslizan de lado; sismos (falla de San Andrés).
\end{definicion}

\begin{ideasclave}
\begin{itemize}
  \item Chile está sobre un borde convergente: la placa de Nazca subduce bajo la Sudamericana.
  \item Magnitud: una por sismo; intensidad: distinta en cada lugar.
  \item Las ondas sísmicas se propagan por el interior de la Tierra.
\end{itemize}
\end{ideasclave}

\begin{metodo}
\begin{enumerate}
  \item Identifica el tipo de borde de placa.
  \item Relaciónalo con el fenómeno (sismo, volcán, relieve).
  \item Distingue epicentro de hipocentro y magnitud de intensidad.
\end{enumerate}
\end{metodo}

\begin{errorfrecuente}[Trampas frecuentes del DEMRE]
\begin{itemize}
  \item Confundir epicentro con hipocentro.
  \item Pensar que la intensidad de un sismo es igual en todas las ciudades.
\end{itemize}
\end{errorfrecuente}

\begin{ejemplo}[Ejemplo resuelto]
{\sffamily\footnotesize\color{zinc500} Enunciado}\par
¿Por qué Chile tiene tantos sismos y volcanes?
\alternativa{A}{Porque está en el centro de una placa tectónica.}
\alternativa{B}{Porque la placa de Nazca subduce bajo la placa Sudamericana.}
\alternativa{C}{Porque las placas se separan bajo el país.}
\alternativa{D}{Porque el núcleo terrestre está más cerca de la superficie en Chile.}
\clave{B}
{\sffamily\footnotesize\color{zinc500} Análisis}\par Chile está en un borde convergente.
{\sffamily\footnotesize\color{zinc500} Desarrollo}
\begin{enumerate}
  \item En la subducción, la placa oceánica se hunde bajo la continental.
  \item La fricción y la fusión del material generan sismos y volcanismo.
\end{enumerate}
\end{ejemplo}

\subsection{Circuitos eléctricos y consumo de energía}

La corriente eléctrica es el flujo de carga por un circuito y se mide en amperes. La ley de Ohm relaciona voltaje, corriente y resistencia: $V=IR$. En serie, la corriente es la misma en todas las resistencias y los voltajes se suman; en paralelo, el voltaje es el mismo y las corrientes se suman.

La potencia eléctrica es $P=VI$ y la energía consumida, $E=Pt$ (en kWh en las cuentas de luz). Un artefacto eficiente transforma en energía útil una mayor parte de la que consume. En una casa, los artefactos se conectan en paralelo, con interruptores, enchufes, fusibles o disyuntores que protegen la instalación y la conexión a tierra.

\begin{ley}[Fórmulas]
\begin{gather*}
V=I\,R \\
P=V\,I \\
E=P\,t \\
R_{\text{serie}}=R_1+R_2 \\
\frac{1}{R_{\text{paralelo}}}=\frac{1}{R_1}+\frac{1}{R_2}
\end{gather*}
\end{ley}

{\sffamily\footnotesize\color{zinc500} Esquema comparativo}
\begin{center}\small
\begin{tabular}{@{}>{\raggedright\arraybackslash}p{\dimexpr(\linewidth-4\tabcolsep)/3\relax}>{\raggedright\arraybackslash}p{\dimexpr(\linewidth-4\tabcolsep)/3\relax}>{\raggedright\arraybackslash}p{\dimexpr(\linewidth-4\tabcolsep)/3\relax}@{}}
\toprule
 & \textbf{Serie} & \textbf{Paralelo} \\
\midrule
Corriente & igual en todas & se reparte \\
Voltaje & se reparte & igual en todas \\
$R$ equivalente & $R_1+R_2$ & menor que cada una \\
\bottomrule
\end{tabular}
\end{center}

\begin{ideasclave}
\begin{itemize}
  \item Las casas conectan los artefactos en paralelo (mismo voltaje).
  \item 1 kWh es la energía de un artefacto de 1000 W funcionando una hora.
  \item La resistencia equivalente en paralelo es menor que la menor de ellas.
\end{itemize}
\end{ideasclave}

\begin{metodo}
\begin{enumerate}
  \item Identifica si las resistencias están en serie o en paralelo.
  \item Calcula la resistencia equivalente y la corriente total.
  \item Aplica Ohm, la potencia o la energía consumida según lo que se pide.
\end{enumerate}
\end{metodo}

\begin{errorfrecuente}[Trampas frecuentes del DEMRE]
\begin{itemize}
  \item Sumar resistencias en paralelo como si estuvieran en serie.
  \item Confundir potencia (W) con energía consumida (kWh).
\end{itemize}
\end{errorfrecuente}

\begin{ejemplo}[Ejemplo resuelto]
{\sffamily\footnotesize\color{zinc500} Enunciado}\par
Un hervidor funciona con 220 V y por él circulan 10 A. ¿Cuál es su potencia?
\alternativa{A}{22 W}
\alternativa{B}{230 W}
\alternativa{C}{2200 W}
\alternativa{D}{22 000 W}
\clave{C}
{\sffamily\footnotesize\color{zinc500} Análisis}\par $P=VI$.
{\sffamily\footnotesize\color{zinc500} Desarrollo}
\begin{enumerate}
  \item $P=220\cdot 10=2200$ W.
\end{enumerate}
\end{ejemplo}

\begin{ejemplo}[Ejemplo tipo PAES 2: Circuitos eléctricos]
{\sffamily\footnotesize\color{zinc500} Enunciado}\par
Dos resistencias de $4\ \Omega$ y $6\ \Omega$ se conectan en serie a una batería de 20 V. ¿Cuál es el voltaje entre los extremos de la resistencia de $6\ \Omega$?
\alternativa{A}{2 V}
\alternativa{B}{8 V}
\alternativa{C}{12 V}
\alternativa{D}{20 V}
\clave{C}
{\sffamily\footnotesize\color{zinc500} Análisis}\par En un circuito en serie la corriente es la misma en ambas resistencias; con $I = 2$ A, el voltaje en la de $6\ \Omega$ es $12$ V.
{\sffamily\footnotesize\color{zinc500} Desarrollo}
\begin{enumerate}
  \item En serie, las resistencias forman un único camino para la corriente, así que la resistencia equivalente es la suma: $R_{\text{eq}} = R_1 + R_2$.
  \item Reemplazando: $R_{\text{eq}} = 4\ \Omega + 6\ \Omega = 10\ \Omega$.
  \item Por la ley de Ohm, $V = I\,R$, así que la corriente total es $I = \dfrac{V}{R_{\text{eq}}} = \dfrac{20\ \mathrm{V}}{10\ \Omega} = 2\ \mathrm{A}$.
  \item En serie, la misma corriente pasa por ambas resistencias, por lo que la de $6\ \Omega$ también es recorrida por $I = 2$ A.
  \item Se aplica la ley de Ohm a esa resistencia: $V_2 = I\,R_2 = 2\ \mathrm{A}\cdot 6\ \Omega = 12\ \mathrm{V}$.
  \item Control: en la de $4\ \Omega$, $V_1 = 2\cdot 4 = 8$ V; y $8 + 12 = 20$ V, que coincide con la batería, como exige la ley de Kirchhoff de voltajes.
  \item La alternativa A (2 V) es la corriente en amperes, no el voltaje; la alternativa B (8 V) corresponde a la resistencia de $4\ \Omega$.
  \item La alternativa D (20 V) es el voltaje total de la batería, que se reparte entre las dos resistencias.
  \item Así, el voltaje en la resistencia de $6\ \Omega$ es de $12$ V. Por lo tanto, la alternativa correcta es la C.
\end{enumerate}
\end{ejemplo}

\section{Química}

{\itshape\color{zinc700} Estructura atómica, química orgánica, reacciones y estequiometría, disoluciones.}

\subsection{Materia y estructura atómica}

La materia se clasifica en sustancias puras (elementos y compuestos) y mezclas, homogéneas o heterogéneas. Las mezclas se separan por métodos físicos: decantación, filtración, tamizado y destilación. Las propiedades físicas, como los puntos de fusión y ebullición o la densidad, permiten identificar sustancias.

Los modelos atómicos evolucionaron con la evidencia: Dalton (esfera indivisible), Thomson (electrones incrustados), Rutherford (núcleo pequeño y denso, por el experimento de la lámina de oro) y Bohr (electrones en niveles de energía). El número atómico $Z$ es la cantidad de protones; el número másico $A$, la suma de protones y neutrones. Los isótopos tienen igual $Z$ y distinto $A$.

\begin{formula}[Relaciones]
\begin{gather*}
n=A-Z \\
e^-=Z-\text{carga} \\
\rho=\frac{m}{V}
\end{gather*}

Modelos: Dalton $\rightarrow$ Thomson $\rightarrow$ Rutherford $\rightarrow$ Bohr.
\end{formula}

\begin{ideasclave}
\begin{itemize}
  \item Un catión perdió electrones; un anión los ganó.
  \item La destilación separa líquidos por su punto de ebullición; la filtración, un sólido de un líquido.
  \item El agua potable es una mezcla; el agua destilada, una sustancia pura.
\end{itemize}
\end{ideasclave}

\begin{metodo}
\begin{enumerate}
  \item Clasifica la muestra (elemento, compuesto o mezcla) o identifica el modelo.
  \item Elige el método de separación según la propiedad que distingue a los componentes.
  \item Calcula protones, neutrones y electrones con $Z$, $A$ y la carga.
\end{enumerate}
\end{metodo}

\begin{errorfrecuente}[Trampas frecuentes del DEMRE]
\begin{itemize}
  \item Tomar $A$ como número de neutrones.
  \item Usar filtración para separar una mezcla homogénea.
\end{itemize}
\end{errorfrecuente}

\begin{ejemplo}[Ejemplo resuelto]
{\sffamily\footnotesize\color{zinc500} Enunciado}\par
El ion ${}^{32}_{16}\mathrm{S}^{2-}$ tiene:
\alternativa{A}{16 protones y 16 electrones.}
\alternativa{B}{16 protones y 18 electrones.}
\alternativa{C}{18 protones y 16 electrones.}
\alternativa{D}{32 protones y 18 electrones.}
\clave{B}
{\sffamily\footnotesize\color{zinc500} Análisis}\par La carga $2-$ indica que ganó dos electrones.
{\sffamily\footnotesize\color{zinc500} Desarrollo}
\begin{enumerate}
  \item Protones: $Z=16$.
  \item Electrones: $16+2=18$.
\end{enumerate}
\end{ejemplo}

\begin{ejemplo}[Ejemplo tipo PAES 2: Estructura atómica]
{\sffamily\footnotesize\color{zinc500} Enunciado}\par
¿Cuántos protones, neutrones y electrones tiene el ion ${}^{23}_{11}\mathrm{Na}^{+}$?
\alternativa{A}{11 protones, 12 neutrones y 11 electrones.}
\alternativa{B}{11 protones, 12 neutrones y 10 electrones.}
\alternativa{C}{12 protones, 11 neutrones y 10 electrones.}
\alternativa{D}{11 protones, 23 neutrones y 10 electrones.}
\alternativa{E}{10 protones, 12 neutrones y 11 electrones.}
\clave{B}
{\sffamily\footnotesize\color{zinc500} Análisis}\par En el átomo neutro, $Z$ es el número de protones; el catión $\mathrm{Na^+}$ perdió un electrón, por lo que tiene 10 electrones y 12 neutrones.
{\sffamily\footnotesize\color{zinc500} Desarrollo}
\begin{enumerate}
  \item El subíndice $Z = 11$ indica el número de protones, que define al elemento sodio.
  \item El superíndice $A = 23$ es el número másico, la suma de protones y neutrones.
  \item Neutrones: $N = A - Z = 23 - 11 = 12$.
  \item El átomo neutro tiene tantos electrones como protones, es decir, 11.
  \item La carga $+1$ indica que perdió un electrón: $11 - 1 = 10$ electrones.
  \item Control: la carga neta es $11 - 10 = +1$, igual a la carga indicada en el símbolo.
  \item La alternativa A corresponde al átomo neutro, pues no considera la pérdida de carga; la C invierte protones y neutrones, con 12 protones, lo que no corresponde a $Z = 11$.
  \item La alternativa D toma el número másico como cantidad de neutrones, y la E cambia el número de protones a 10, lo que correspondería a otro elemento.
  \item Así, el ion tiene 11 protones, 12 neutrones y 10 electrones. Por lo tanto, la alternativa correcta es la B.
\end{enumerate}
\end{ejemplo}

\begin{ejemplo}[Ejemplo tipo PAES 3: Configuración electrónica y tabla periódica]
{\sffamily\footnotesize\color{zinc500} Enunciado}\par
Un elemento tiene número atómico $Z = 17$. ¿Cuál de las siguientes opciones describe correctamente su ubicación en la tabla periódica y su comportamiento químico?
\alternativa{A}{Período 3, grupo 17; tiende a ganar un electrón.}
\alternativa{B}{Período 3, grupo 7; tiende a perder siete electrones.}
\alternativa{C}{Período 2, grupo 17; tiende a ganar un electrón.}
\alternativa{D}{Período 3, grupo 1; tiende a perder un electrón.}
\alternativa{E}{Período 4, grupo 17; tiende a compartir un electrón.}
\clave{A}
{\sffamily\footnotesize\color{zinc500} Análisis}\par Con $Z = 17$ la configuración termina en $3p^5$: el elemento está en el período 3, en el grupo 17, con 7 electrones de valencia, y tiende a ganar un electrón.
{\sffamily\footnotesize\color{zinc500} Desarrollo}
\begin{enumerate}
  \item Con $Z = 17$, el átomo neutro tiene 17 electrones, distribuidos según el principio de Aufbau: $1s^2\,2s^2\,2p^6\,3s^2\,3p^5$.
  \item Se comprueba la suma: $2 + 2 + 6 + 2 + 5 = 17$ electrones.
  \item El nivel más externo es $n = 3$, así que el elemento pertenece al período 3.
  \item Los electrones de valencia son los del nivel $n = 3$: $3s^2\,3p^5$, es decir, $2 + 5 = 7$ electrones.
  \item Siete electrones de valencia corresponden al grupo 17 (halógenos) en la numeración de 18 grupos.
  \item Le falta un electrón para completar el octeto, por lo que tiende a ganarlo y formar el anión $\mathrm{Cl^-}$, con configuración de gas noble.
  \item La alternativa B confunde el número de grupo con la tendencia y propone perder siete electrones, algo energéticamente inviable; la C indica período 2, cuando el nivel externo es 3.
  \item La alternativa D indica grupo 1 y perder un electrón, propio de los metales alcalinos; la E indica período 4, y compartir un electrón no es la tendencia típica de este grupo.
  \item Así, el elemento está en el período 3, grupo 17, y tiende a ganar un electrón. Por lo tanto, la alternativa correcta es la A.
\end{enumerate}
\end{ejemplo}

\subsection{Química orgánica}

La química orgánica estudia los compuestos del carbono. El carbono forma cuatro enlaces y puede unirse a otros carbonos formando cadenas y anillos. Su hibridación determina la geometría: $sp^3$ (enlaces simples, tetraédrica), $sp^2$ (un doble enlace, plana) y $sp$ (triple enlace, lineal).

Los grupos funcionales determinan las propiedades: alcoholes ($-\mathrm{OH}$), aldehídos ($-\mathrm{CHO}$), cetonas ($-\mathrm{CO}-$), ácidos carboxílicos ($-\mathrm{COOH}$), éteres ($-\mathrm{O}-$), ésteres, aminas y amidas. Los isómeros tienen igual fórmula molecular y distinta estructura.

\begin{definicion}[Grupos funcionales]
Alcohol: $\mathrm{R{-}OH}$. Aldehído: $\mathrm{R{-}CHO}$. Cetona: $\mathrm{R{-}CO{-}R'}$. Ácido: $\mathrm{R{-}COOH}$. Éter: $\mathrm{R{-}O{-}R'}$. Éster: $\mathrm{R{-}COO{-}R'}$. Amina: $\mathrm{R{-}NH_2}$.
\end{definicion}

\begin{ideasclave}
\begin{itemize}
  \item Cada doble enlace aporta un enlace pi ($\pi$).
  \item Carbono con doble enlace: $sp^2$.
  \item Alcanos: solo enlaces simples ($\mathrm{C_nH_{2n+2}}$).
\end{itemize}
\end{ideasclave}

\begin{metodo}
\begin{enumerate}
  \item Identifica el grupo funcional o los tipos de enlace.
  \item Cuenta carbonos y enlaces múltiples.
  \item Nombra o clasifica según las reglas IUPAC.
\end{enumerate}
\end{metodo}

\begin{errorfrecuente}[Trampas frecuentes del DEMRE]
\begin{itemize}
  \item Confundir aldehído (carbonilo en el extremo) con cetona (carbonilo interno).
  \item Contar mal los enlaces pi.
\end{itemize}
\end{errorfrecuente}

\begin{ejemplo}[Ejemplo resuelto]
{\sffamily\footnotesize\color{zinc500} Enunciado}\par
¿Cuál es la fórmula molecular del propano, un alcano de 3 carbonos?
\alternativa{A}{$\mathrm{C_3H_4}$}
\alternativa{B}{$\mathrm{C_3H_6}$}
\alternativa{C}{$\mathrm{C_3H_8}$}
\alternativa{D}{$\mathrm{C_3H_{10}}$}
\clave{C}
{\sffamily\footnotesize\color{zinc500} Análisis}\par Alcanos: $\mathrm{C_nH_{2n+2}}$.
{\sffamily\footnotesize\color{zinc500} Desarrollo}
\begin{enumerate}
  \item $n=3$: $2\cdot 3+2=8$ hidrógenos.
  \item $\mathrm{C_3H_8}$.
\end{enumerate}
\end{ejemplo}

\begin{ejemplo}[Ejemplo tipo PAES 2: Estequiometría]
{\sffamily\footnotesize\color{zinc500} Enunciado}\par
Considera la reacción balanceada:

\[2\,\mathrm{H_2} + \mathrm{O_2} \rightarrow 2\,\mathrm{H_2O}\]

Si reaccionan por completo 4 g de $\mathrm{H_2}$ con la cantidad exacta de $\mathrm{O_2}$, ¿qué masa de agua se forma? (Masas molares: H = 1 g/mol; O = 16 g/mol.)
\alternativa{A}{18 g}
\alternativa{B}{34 g}
\alternativa{C}{36 g}
\alternativa{D}{72 g}
\clave{C}
{\sffamily\footnotesize\color{zinc500} Análisis}\par La estequiometría convierte 4 g de hidrógeno en 2 mol de $\mathrm{H_2}$, que producen 2 mol de agua, es decir, 36 g de $\mathrm{H_2O}$.
{\sffamily\footnotesize\color{zinc500} Desarrollo}
\begin{enumerate}
  \item Primero se calcula la masa molar del agua: $M(\mathrm{H_2O}) = 2\cdot 1 + 16 = 18$ g/mol.
  \item Se convierten los 4 g de $\mathrm{H_2}$ a moles, con $M(\mathrm{H_2}) = 2$ g/mol: $n = \dfrac{4\ \text{g}}{2\ \text{g/mol}} = 2$ mol.
  \item La ecuación indica la proporción 2 mol de $\mathrm{H_2}$ : 2 mol de $\mathrm{H_2O}$, por lo que se forman 2 mol de agua.
  \item Se convierte a masa: $m = n\cdot M = 2\ \text{mol}\cdot 18\ \text{g/mol} = 36$ g.
  \item Control por conservación de la masa: los reactivos suman 4 g de $\mathrm{H_2}$ y 32 g de $\mathrm{O_2}$ (que son 1 mol según la ecuación), es decir, 36 g, igual a la masa de agua formada.
  \item La alternativa B (34 g) suma 2 g de hidrógeno (1 mol) a 32 g de oxígeno, como si no se usaran los coeficientes de la ecuación.
  \item La alternativa A (18 g) es la masa de un mol de agua, pues no se consideró la cantidad de moles formados.
  \item La alternativa D (72 g) usa 4 mol de agua, lo que confunde la masa de hidrógeno en gramos con su cantidad en moles.
  \item Así, la masa de agua formada es de 36 g. Por lo tanto, la alternativa correcta es la C.
\end{enumerate}
\end{ejemplo}

\begin{ejemplo}[Ejemplo tipo PAES 3: pH]
{\sffamily\footnotesize\color{zinc500} Enunciado}\par
Una solución tiene una concentración de iones hidrógeno $[\mathrm{H^+}] = 1\cdot 10^{-3}$ mol/L a 25 \textdegree{}C. ¿Cuál es su pH y cómo se clasifica?
\alternativa{A}{pH = 11; básica}
\alternativa{B}{pH = 3; básica}
\alternativa{C}{pH = $-$3; ácida}
\alternativa{D}{pH = 3; ácida}
\clave{D}
{\sffamily\footnotesize\color{zinc500} Análisis}\par Con $[\mathrm{H^+}] = 10^{-3}$ mol/L, el pH es 3; a 25 \textdegree{}C, un pH menor que 7 corresponde a una solución ácida.
{\sffamily\footnotesize\color{zinc500} Desarrollo}
\begin{enumerate}
  \item El pH se define como $\mathrm{pH} = -\log[\mathrm{H^+}]$, con la concentración expresada en mol/L.
  \item Se reemplaza: $\mathrm{pH} = -\log(1\cdot 10^{-3})$.
  \item Como $\log(10^{-3}) = -3$, resulta $\mathrm{pH} = -(-3) = 3$.
  \item Control: con $\mathrm{pH} = 3$ el $\mathrm{pOH}$ es $14 - 3 = 11$, lo que implica $[\mathrm{OH^-}] = 10^{-11}$ mol/L, consistente con una solución muy ácida.
  \item Clasificación a 25 \textdegree{}C: pH 7 es neutro, pH menor que 7 es ácido y pH mayor que 7 es básico. Como $3 < 7$, la solución es ácida.
  \item La alternativa A (pH 11, básica) es el pOH calculado como si fuera pH.
  \item La alternativa B (pH 3, básica) acierta el valor, pero invierte la clasificación: un pH bajo indica acidez.
  \item La alternativa C (pH $-$3) olvida el signo negativo de la definición, lo que da un valor que no corresponde a esta concentración.
  \item Así, la solución tiene pH 3 y es ácida. Por lo tanto, la alternativa correcta es la D.
\end{enumerate}
\end{ejemplo}

\subsection{Reacciones químicas y estequiometría}

En una reacción química los átomos se reordenan: la masa se conserva (ley de Lavoisier), por eso la ecuación debe estar balanceada. Los coeficientes indican la proporción en moles.

El mol es la cantidad de sustancia que contiene $6{,}02\cdot 10^{23}$ partículas. La masa molar convierte gramos en moles. El reactivo limitante es el que se consume primero y determina la cantidad de producto.

\begin{formula}[Relaciones]
\begin{gather*}
n=\frac{m}{M} \\
N=n\cdot 6{,}02\cdot 10^{23} \\
\%\,\text{elemento}=\frac{\text{masa del elemento}}{\text{masa molar}}\cdot 100
\end{gather*}

Leyes ponderales: conservación de la masa (Lavoisier), proporciones definidas (Proust) y múltiples (Dalton). La fórmula empírica es la proporción mínima entre átomos; la molecular, un múltiplo de ella.
\end{formula}

\begin{ideasclave}
\begin{itemize}
  \item Balancea antes de calcular.
  \item Los coeficientes relacionan moles, no gramos.
  \item Masa total de reactantes = masa total de productos.
\end{itemize}
\end{ideasclave}

\begin{metodo}
\begin{enumerate}
  \item Balancea la ecuación.
  \item Convierte los datos a moles.
  \item Usa la proporción de coeficientes y vuelve a la unidad pedida.
\end{enumerate}
\end{metodo}

\begin{errorfrecuente}[Trampas frecuentes del DEMRE]
\begin{itemize}
  \item Usar las proporciones de coeficientes directamente con gramos.
  \item Olvidar balancear la ecuación.
\end{itemize}
\end{errorfrecuente}

\begin{ejemplo}[Ejemplo resuelto]
{\sffamily\footnotesize\color{zinc500} Enunciado}\par
En la reacción balanceada $\mathrm{N_2}+3\,\mathrm{H_2}\rightarrow 2\,\mathrm{NH_3}$, ¿cuántos moles de $\mathrm{NH_3}$ se forman a partir de 6 mol de $\mathrm{H_2}$ con exceso de $\mathrm{N_2}$?
\alternativa{A}{2 mol}
\alternativa{B}{3 mol}
\alternativa{C}{4 mol}
\alternativa{D}{6 mol}
\clave{C}
{\sffamily\footnotesize\color{zinc500} Análisis}\par La proporción entre $\mathrm{H_2}$ y $\mathrm{NH_3}$ es $3:2$.
{\sffamily\footnotesize\color{zinc500} Desarrollo}
\begin{enumerate}
  \item $6\ \text{mol H}_2\cdot\dfrac{2\ \text{mol NH}_3}{3\ \text{mol H}_2}=4$ mol de $\mathrm{NH_3}$.
\end{enumerate}
\end{ejemplo}

\subsection{Disoluciones y concentración}

Una disolución es una mezcla homogénea de soluto y disolvente. Su concentración se expresa en unidades físicas (\% m/m, \% m/v, \% v/v y ppm) o químicas (molaridad, molalidad y fracción molar). Al diluir, la cantidad de soluto se mantiene; al mezclar dos disoluciones del mismo soluto, se suman los solutos y los volúmenes.

La solubilidad es la máxima cantidad de soluto que se disuelve en una cantidad de disolvente a cierta temperatura; depende de la naturaleza de las sustancias, la temperatura y, en los gases, la presión.

\begin{formula}[Fórmulas]
\begin{gather*}
M=\frac{n_{\text{soluto}}}{V_{\text{disolución}}\,(\text{L})} \\
m=\frac{n_{\text{soluto}}}{\text{kg disolvente}} \\
\%\,m/v=\frac{g\ \text{soluto}}{mL\ \text{disolución}}\cdot 100 \\
C_1V_1=C_2V_2
\end{gather*}
\end{formula}

\begin{ideasclave}
\begin{itemize}
  \item Diluir no cambia los moles de soluto.
  \item Molaridad: por litro de \textbf{disolución}; molalidad: por kilogramo de \textbf{disolvente}.
  \item La solubilidad de los gases baja al subir la temperatura.
\end{itemize}
\end{ideasclave}

\begin{metodo}
\begin{enumerate}
  \item Convierte los datos a moles, litros o gramos según la unidad pedida.
  \item Aplica la fórmula de concentración, de dilución o de mezcla.
  \item Revisa que el resultado tenga sentido (diluir siempre baja la concentración).
\end{enumerate}
\end{metodo}

\begin{errorfrecuente}[Trampas frecuentes del DEMRE]
\begin{itemize}
  \item Usar mililitros en vez de litros en la molaridad.
  \item Confundir masa de disolvente con masa de disolución.
\end{itemize}
\end{errorfrecuente}

\begin{ejemplo}[Ejemplo resuelto]
{\sffamily\footnotesize\color{zinc500} Enunciado}\par
Se diluyen 100 mL de una disolución 2 mol/L hasta 400 mL. ¿Cuál es la nueva concentración?
\alternativa{A}{0,5 mol/L}
\alternativa{B}{1 mol/L}
\alternativa{C}{4 mol/L}
\alternativa{D}{8 mol/L}
\clave{A}
{\sffamily\footnotesize\color{zinc500} Análisis}\par Los moles de soluto se conservan: $C_1V_1=C_2V_2$.
{\sffamily\footnotesize\color{zinc500} Desarrollo}
\begin{enumerate}
  \item $2\cdot 100=C_2\cdot 400$.
  \item $C_2=0{,}5$ mol/L.
\end{enumerate}
\end{ejemplo}

\begin{ejemplo}[Ejemplo tipo PAES 2: Concentración molar]
{\sffamily\footnotesize\color{zinc500} Enunciado}\par
Se disuelven 20 g de hidróxido de sodio ($\mathrm{NaOH}$, masa molar 40 g/mol) en agua hasta completar 500 mL de solución. ¿Cuál es la concentración molar de la solución?
\alternativa{A}{0,04 mol/L}
\alternativa{B}{0,5 mol/L}
\alternativa{C}{1 mol/L}
\alternativa{D}{2 mol/L}
\alternativa{E}{40 mol/L}
\clave{C}
{\sffamily\footnotesize\color{zinc500} Análisis}\par La molaridad es moles de soluto sobre litros de solución: 20 g de NaOH equivalen a 0,5 mol disueltos en 0,5 L, es decir, 1 mol/L.
{\sffamily\footnotesize\color{zinc500} Desarrollo}
\begin{enumerate}
  \item La concentración molar se define como $M = \dfrac{n}{V}$, con $n$ en mol y $V$ en litros.
  \item Se calcula la cantidad de soluto: $n = \dfrac{20\ \text{g}}{40\ \text{g/mol}} = 0{,}5$ mol.
  \item Se convierte el volumen a litros: $500\ \text{mL}\cdot \dfrac{1\ \text{L}}{1000\ \text{mL}} = 0{,}5$ L.
  \item Se reemplaza en la definición: $M = \dfrac{0{,}5\ \text{mol}}{0{,}5\ \text{L}}$.
  \item Se simplifica: $M = 1$ mol/L.
  \item Control: una solución de 1 mol/L contiene 1 mol en cada litro; en medio litro habría 0,5 mol, que es justo lo que se disolvió.
  \item La alternativa B (0,5 mol/L) entrega los moles de soluto sin dividir por el volumen; la alternativa A (0,04 mol/L) divide gramos entre mililitros, $20/500$.
  \item La alternativa D (2 mol/L) corresponde a usar 0,25 L como volumen, un error de conversión; la alternativa E (40 mol/L) divide $20$ g entre $0{,}5$ L sin pasar los gramos a moles.
  \item Así, la concentración molar es de 1 mol/L. Por lo tanto, la alternativa correcta es la C.
\end{enumerate}
\end{ejemplo}

\section{Habilidades científicas}

{\itshape\color{zinc700} Pensamiento científico: variables, hipótesis, procedimientos y conclusiones.}

\subsection{Habilidades de pensamiento científico}

Muchas preguntas de Ciencias no evalúan un contenido, sino el razonamiento experimental. La variable independiente es la que el investigador manipula; la dependiente, la que mide como respuesta; las controladas, las que se mantienen constantes.

Una hipótesis es una respuesta tentativa y comprobable a una pregunta de investigación. Una conclusión se deriva de los resultados obtenidos y no debe ir más allá de ellos. Una teoría explica un conjunto amplio de fenómenos; una ley describe una regularidad, a menudo matemática.

\begin{mecanismo}[Esquema de una investigación]
Observación $\rightarrow$ pregunta de investigación $\rightarrow$ hipótesis $\rightarrow$ diseño experimental (variables y grupo control) $\rightarrow$ resultados $\rightarrow$ conclusión.
\end{mecanismo}

\begin{ideasclave}
\begin{itemize}
  \item Independiente: lo que yo cambio. Dependiente: lo que mido.
  \item Una conclusión solo puede afirmar lo que los datos muestran.
  \item El grupo control permite comparar el efecto de la variable.
\end{itemize}
\end{ideasclave}

\begin{metodo}
\begin{enumerate}
  \item Identifica qué se cambia a propósito y qué se mide.
  \item Revisa qué factores se mantienen constantes.
  \item Evalúa si la conclusión se sigue de los datos, sin generalizar.
\end{enumerate}
\end{metodo}

\begin{errorfrecuente}[Trampas frecuentes del DEMRE]
\begin{itemize}
  \item Confundir la variable que se mide con la que se manipula.
  \item Elegir una conclusión que generaliza más allá de los datos.
\end{itemize}
\end{errorfrecuente}

\begin{ejemplo}[Ejemplo resuelto]
{\sffamily\footnotesize\color{zinc500} Enunciado}\par
Una estudiante riega cuatro plantas iguales con distinta cantidad de agua y mide su altura después de un mes. ¿Cuál es la variable dependiente?
\alternativa{A}{La cantidad de agua.}
\alternativa{B}{El tipo de planta.}
\alternativa{C}{La altura de las plantas.}
\alternativa{D}{El tiempo de un mes.}
\clave{C}
{\sffamily\footnotesize\color{zinc500} Análisis}\par La dependiente es la que se mide como respuesta.
{\sffamily\footnotesize\color{zinc500} Desarrollo}
\begin{enumerate}
  \item Manipula: cantidad de agua (independiente).
  \item Mide: altura (dependiente). Tipo de planta y tiempo son controladas.
\end{enumerate}
\end{ejemplo}
%%</CONTENIDO>%%

\end{document}
